% 第一章 绪论
% ch1-introduction.tex

\subsection{研究背景}

奥林匹克运动会作为全球最具影响力的体育盛事，其项目的入选与淘汰直接影响着全球体育运动的发展方向。近年来，奥运项目的评估与调整日益频繁，如2020年东京奥运会新增攀岩、冲浪、滑板等项目，2024年巴黎奥运会引入霹雳舞，而2028年洛杉矶奥运会又将取消部分项目。这些决策的背后，反映出国际奥委会（IOC）对项目评估标准的不断探索与调整。

然而，当前的奥运项目评估主要依赖专家委员会的主观判断，缺乏系统化、可量化的决策工具。这种传统方法存在以下问题：
\begin{enumerate}
\item 主观性强：专家意见存在个体差异，难以达成共识
\item 缺乏量化依据：评估标准模糊，决策透明度不足
\item 忽视客观信息：未充分利用历史数据和统计信息
\end{enumerate}

因此，构建一个科学、客观、可量化的奥运项目评估模型，对于提升IOC决策的科学性和透明度具有重要意义。

\subsection{研究意义}

\subsubsection{理论意义}

本研究将多准则决策方法（MCDM）应用于奥运项目评估领域，创新性地采用AHP-EWM混合权重模型，为体育决策科学化提供了新的研究范式。通过建立六维评价指标体系，丰富了奥运项目管理理论，为后续研究提供了可借鉴的方法论框架。

\subsubsection{实践意义}

本研究的成果可直接应用于：
\begin{enumerate}
\item 为IOC提供定量化的项目评估工具
\item 为各国奥委会的项目申办提供决策支持
\item 为国际体育联合会的项目推广提供策略建议
\item 为体育管理教育提供案例研究素材
\end{enumerate}

\subsection{国内外研究现状}

\subsubsection{多准则决策方法研究}

多准则决策（Multi-Criteria Decision Making, MCDM）是决策科学的重要分支。Saaty\cite{saaty2007analytic}于1970年代提出层次分析法（AHP），通过构建判断矩阵量化专家的主观判断。熵权法（EWM）作为一种客观赋权方法，通过信息熵度量指标的信息量，为多准则决策提供了客观依据。

近年来，AHP与熵权法（EWM）的组合应用成为研究热点。Wang和Liu\cite{wang2019comprehensive}在海洋软实力评估中验证了EWM的有效性；Zhang等\cite{zhang2023risk}将AHP-EWM应用于灾害风险评估，证明了混合模型的优势。

\subsubsection{奥运项目评估研究}

Song等\cite{song2025quantifying}在2025年提出了基于AHP-PCA-KNN的奥运项目选择框架，系统总结了IOC的六维评估标准，为本研究提供了重要参考。李凤梅\cite{li2013xiaji}从人文、经济、社会三维度分析了奥运项目设置标准，为指标体系设计提供了思路。

\subsection{研究内容与方法}

\subsubsection{研究内容}

本研究主要包括以下内容：
\begin{enumerate}
\item 构建六维评价指标体系
\item 设计AHP-EWM混合权重计算方法
\item 收集历史数据并验证模型
\item 开发评估系统与可视化仪表盘
\item 对2032年奥运项目进行预测分析
\end{enumerate}

\subsubsection{研究方法}

本研究采用以下方法：
\begin{enumerate}
\item 文献研究法：梳理AHP、EWM及奥运评估相关文献
\item 专家访谈法：收集专家对指标重要性的判断
\item 实证分析法：用历史数据验证模型有效性
\item 系统开发法：实现评估系统原型
\end{enumerate}

\subsection{论文组织结构}

本文共分六章：
\begin{enumerate}
\item 第一章：绪论，介绍研究背景、意义及主要内容
\item 第二章：理论基础与文献综述，阐述相关理论与研究进展
\item 第三章：混合权重评估模型构建，详细推导数学模型
\item 第四章：实验设计与数据分析，验证模型有效性
\item 第五章：系统设计与实现，展示评估系统
\item 第六章：总结与展望，总结研究贡献并指出未来方向
\end{enumerate}